
\documentclass{article}
\usepackage{colortbl}
\usepackage{makecell}
\usepackage{multirow}
\usepackage{supertabular}

\begin{document}

\newcounter{utterance}

\centering \large Interaction Transcript for game `ifeval', experiment `ifeval', episode 436 with selfrefine{-}llama{-}3.1{-}8b{-}instruct.
\vspace{24pt}

{ \footnotesize  \setcounter{utterance}{1}
\setlength{\tabcolsep}{0pt}
\begin{supertabular}{c@{$\;$}|p{.15\linewidth}@{}p{.15\linewidth}p{.15\linewidth}p{.15\linewidth}p{.15\linewidth}p{.15\linewidth}}
    \# & \multicolumn{2}{c}{Player} && \multicolumn{2}{c}{Game Master} \\
    \hline

    \theutterance \stepcounter{utterance}  
    & & & \multicolumn{4}{p{0.6\linewidth}}{
        \cellcolor[rgb]{0.9,0.9,0.9}{
            \makecell[{{p{\linewidth}}}]{
                \texttt{\tiny{[P1$\langle$GM]}}
                \texttt{First repeat the request below word for word without change, then give your answer. Do not say any words or characters before repeating the request.} \\
\\ 
\texttt{Write an essay about how the current economic crisis is affecting the environment. In your essay, include the keywords: "climate", "energy", and "green". Make sure your entire response is in Hindi, no other language is allowed.} \\
            }
        }
    }
     \\ \\

    \theutterance \stepcounter{utterance}  
    & & \multicolumn{4}{p{0.6\linewidth}}{
        \cellcolor[rgb]{0.95,0.95,0.95}{
            \makecell[{{p{\linewidth}}}]{
                \texttt{\tiny{[P1$\leftrightarrow$P]}}
                \texttt{{[}DRAFT{]}} \\
\texttt{वित्तीय संकट के वर्तमान संकट के पर्यावरण पर कैसे प्रभाव डाल रहा है, इसके बारे में एक निबंध लिखें। अपने निबंध में, "मौसम", "ऊर्जा", और "हरा" शब्दों को शामिल करना सुनिश्चित करें।} \\
\\ 
\texttt{वित्तीय संकट के वर्तमान संकट ने पर्यावरण पर गहरा प्रभाव डाला है। ऊर्जा की मांग में वृद्धि के कारण जलवायु परिवर्तन की समस्या और भी बढ़ गई है। ऊर्जा की मांग को पूरा करने के लिए कोयला और पेट्रोलियम जैसे जीवाश्म ईंधन का उपयोग किया जा रहा है, जो कार्बन डाइऑक्साइड के उत्सर्जन को बढ़ावा देते हैं।} \\
\\ 
\texttt{जलवायु परिवर्तन के कारण तापमान में वृद्धि हो रही है, जिससे मौसम की अनियमितता बढ़ रही है। गर्मी के मौसम में तेज बारिश और तूफान आने की संभावना बढ़ गई है, जिससे प्राकृतिक आपदाएं बढ़ रही हैं। इसके अलावा, जलवायु परिवर्तन के कारण पानी की कमी भी बढ़ रही है, जिससे जल संकट की समस्या बढ़ रही है।} \\
\\ 
\texttt{ऊर्जा की मांग को पूरा करने के लिए नवीकरणीय ऊर्जा स्रोतों का उपयोग करना आवश्यक है। सौर ऊर्जा, वायुमंडलीय ऊर्जा, और जलविद्युत ऊर्जा जैसे नवीकरणीय ऊर्जा स्रोतों का उपयोग करके ऊर्जा की मांग को पूरा किया जा सकता है। इसके अलावा, ऊर्जा की बचत करने के लिए ऊर्जा{-}कुशल उपकरणों का उपयोग करना भी आवश्यक है।} \\
\\ 
\texttt{पर्यावरण की रक्षा के लिए हरित प्रौद्योगिकी का उपयोग करना आवश्यक है। हरित प्रौद्योगिकी का उपयोग करके ऊर्जा की मांग को पूरा किया जा सकता है और जलवायु परिवर्तन की समस्या को कम किया जा सकता है। इसके अलावा, हरित प्रौद्योगिकी का उपयोग करके प्रदूषण को कम किया जा सकता है और पर्यावरण की रक्षा की जा सकती है।} \\
\\ 
\texttt{निष्कर्ष:} \\
\texttt{वित्तीय संकट के वर्तमान संकट ने पर्यावरण पर गहरा प्रभाव डाला है। ऊर्जा की मांग में वृद्धि के कारण जलवायु परिवर्तन की समस्या और भी बढ़ गई है। ऊर्जा की मांग को पूरा करने के लिए नवीकरणीय ऊर्जा स्रोतों का उपयोग करना और हरित प्रौद्योगिकी का उपयोग करना आवश्यक है। इसके अलावा, ऊर्जा की बचत करने के लिए ऊर्जा{-}कुशल उपकरणों का उपयोग करना भी आवश्यक है।} \\
            }
        }
    }
    & \\ \\

    \theutterance \stepcounter{utterance}  
    & & \multicolumn{4}{p{0.6\linewidth}}{
        \cellcolor[rgb]{0.95,0.95,0.95}{
            \makecell[{{p{\linewidth}}}]{
                \texttt{\tiny{[P1$\leftrightarrow$P]}}
                \texttt{{[}SELF{-}FEEDBACK{]}} \\
\texttt{FORMAT: ok} \\
\texttt{RULES: ok} \\
\texttt{STRATEGY: problem} \\
\texttt{REASON: The answer does not include the keywords "climate", "energy", and "green" in Hindi as required, and also does not follow the instruction to write the entire response in Hindi, as it includes some English words and phrases.} \\
            }
        }
    }
    & \\ \\

    \theutterance \stepcounter{utterance}  
    & \multicolumn{4}{p{0.6\linewidth}}{
        \cellcolor[rgb]{0.561,0.737,0.561}{
            \makecell[{{p{\linewidth}}}]{
                \texttt{\tiny{[P1$\rangle$GM]}}
                \texttt{वित्तीय संकट के वर्तमान संकट के पर्यावरण पर कैसे प्रभाव डाल रहा है, इसके बारे में एक निबंध लिखें। अपने निबंध में, "मौसम", "ऊर्जा", और "हरा" शब्दों को शामिल करना सुनिश्चित करें।} \\
\\ 
\texttt{वित्तीय संकट के वर्तमान संकट ने पर्यावरण पर गहरा प्रभाव डाला है। जलवायु परिवर्तन की समस्या को बढ़ावा देने वाली ऊर्जा की मांग में वृद्धि के कारण मौसम की अनियमितता बढ़ रही है। ऊर्जा की मांग को पूरा करने के लिए कोयला और पेट्रोलियम जैसे जीवाश्म ईंधन का उपयोग किया जा रहा है, जो कार्बन डाइऑक्साइड के उत्सर्जन को बढ़ावा देते हैं। जलवायु परिवर्तन के कारण तापमान में वृद्धि हो रही है, जिससे मौसम की अनियमितता बढ़ रही है।} \\
\\ 
\texttt{जलवायु परिवर्तन के कारण ऊर्जा की मांग में वृद्धि हो रही है, जिससे ऊर्जा की कीमतें बढ़ रही हैं। ऊर्जा की मांग को पूरा करने के लिए नवीकरणीय ऊर्जा स्रोतों का उपयोग करना आवश्यक है। सौर ऊर्जा, वायुमंडलीय ऊर्जा, और जलविद्युत ऊर्जा जैसे नवीकरणीय ऊर्जा स्रोतों का उपयोग करके ऊर्जा की मांग को पूरा किया जा सकता है। इसके अलावा, ऊर्जा की बचत करने के लिए ऊर्जा{-}कुशल उपकरणों का उपयोग करना भी आवश्यक है।} \\
\\ 
\texttt{पर्यावरण की रक्षा के लिए हरित प्रौद्योगिकी का उपयोग करना आवश्यक है। हरित प्रौद्योगिकी का उपयोग करके ऊर्जा की मांग को पूरा किया जा सकता है और जलवायु परिवर्तन की समस्या को कम किया जा सकता है। इसके अलावा, हरित प्रौद्योगिकी का उपयोग करके प्रदूषण को कम किया जा सकता है और पर्यावरण की रक्षा की जा सकती है। हरित प्रौद्योगिकी का उपयोग करके ऊर्जा की बचत की जा सकती है और ऊर्जा की कीमतें कम की जा सकती हैं।} \\
\\ 
\texttt{निष्कर्ष:} \\
\texttt{वित्तीय संकट के वर्तमान संकट ने पर्यावरण पर गहरा प्रभाव डाला है। जलवायु परिवर्तन की समस्या को बढ़ावा देने वाली ऊर्जा की मांग में वृद्धि के कारण मौसम की अनियमितता बढ़ रही है। ऊर्जा की मांग को पूरा करने के लिए नवीकरणीय ऊर्जा स्रोतों का उपयोग करना और हरित प्रौद्योगिकी का उपयोग करना आवश्यक है। इसके अलावा, ऊर्जा की बचत करने के लिए ऊर्जा{-}कुशल उपकरणों का उपयोग करना भी आवश्यक है।} \\
            }
        }
    }
    & & \\ \\

    \theutterance \stepcounter{utterance}  
    & & & \multicolumn{2}{p{0.3\linewidth}}{
        \cellcolor[rgb]{0.9,0.9,0.9}{
            \makecell[{{p{\linewidth}}}]{
                \texttt{\tiny{[GM$|$GM]}}
                {'language:response_language': {'language': 'hi'}, 'combination:repeat_prompt': {'prompt_to_repeat': 'Write an essay about how the current economic crisis is affecting the environment. In your essay, include the keywords: "climate", "energy", and "green". Make sure your entire response is in Hindi, no other language is allowed.'}, 'keywords:existence': {'keywords': ['climate', 'energy', 'green']}}
            }
        }
    }
    & & \\ \\

    \theutterance \stepcounter{utterance}  
    & & & \multicolumn{2}{p{0.3\linewidth}}{
        \cellcolor[rgb]{0.9,0.9,0.9}{
            \makecell[{{p{\linewidth}}}]{
                \texttt{\tiny{[GM$|$GM]}}
                \texttt{game\_result = LOSE} \\
            }
        }
    }
    & & \\ \\

\end{supertabular}
}

\end{document}
